\chapter{Retained-Teacher Neural OT Warm Starts}
\label{ch:bf-dpf-learned-ot}

The previous chapter fixed the entropic OT and Sinkhorn layer as a fully declared
teacher object.  The present chapter introduces the narrowest neural extension of
that teacher: do not replace the teacher; learn how to start it better.

This is the main implementation-bearing neural chapter of the OT block.  Its
question is narrower than ``can a network learn transport?''  It asks instead:
can a network predict a good internal solver state so that a retained corrective
Sinkhorn phase needs less work while still protecting the declared teacher object?

\section{Objects and conventions}
\label{sec:bf-learned-ot-objects}

Let \(X=(x_1,\ldots,x_N)^\top\in\mathbb R^{N\times d_x}\) be the weighted source
cloud and \(w\in\Delta_N\) its normalized weights.  Let
\(u=(1/N,\ldots,1/N)\) be the equal target marginal, and let
\(C(X,Z)\in\mathbb R^{N\times N}\) be the declared transport cost matrix.  The
teacher coupling is either the exact regularized OT optimizer
\(\Pi_\varepsilon^\star\) or the finite numerical teacher
\(\Pi_{\varepsilon,K}\).  The barycentric teacher output is
\begin{equation}
\label{eq:bf-learned-ot-barycentric-teacher}
  Y_\tau(X,w,\varepsilon)=B(\Pi_\tau(X,w,\varepsilon)).
\end{equation}

The student does not directly define a new target cloud at first.  It predicts a
solver-relevant latent quantity such as dual variables or scaling vectors.  Write
\begin{equation}
\label{eq:bf-learned-ot-student-map}
  h_\psi = S_\psi(X,w;\varepsilon,N,d_x),
\end{equation}
where \(h_\psi\) is later translated into the initialization of the retained
corrective Sinkhorn solver.

\section{What the teacher computes}
\label{sec:bf-learned-ot-hierarchy}

There is no single teacher unless the variant is named.  In this chapter, the
main teacher is the retained EOT/Sinkhorn route from the previous chapter.  The
most important teacher objects are:
\begin{enumerate}[label=(\roman*)]
  \item the exact regularized coupling \(\Pi_\varepsilon^\star\),
  \item the finite numerical teacher \(\Pi_{\varepsilon,K}\),
  \item the barycentric teacher output \(Y_\tau=B(\Pi_\tau)\).
\end{enumerate}
The chapter’s neural route is only intended to accelerate the finite teacher, not
to redefine these objects.

\section{Algorithm: retained-teacher warm-started Sinkhorn}
\label{sec:bf-learned-ot-teacher}

\begin{algorithm}[ht]
\caption{\texttt{train\_and\_deploy\_warmstarted\_sinkhorn}}
\begin{algorithmic}[1]
\Require training distribution \(\mathcal D\) over weighted particle clouds, declared cost construction, regularization schedule \(\varepsilon\), teacher Sinkhorn solver, architecture class \(S_\psi\), deployment-time corrective budget \(K_{\mathrm{corr}}\)
\Ensure trained predictor \(S_\psi\), corrected deployment coupling \(\widehat\Pi_{\psi,\varepsilon,K_{\mathrm{corr}}}\), barycentric output cloud \(\widetilde X_{\psi,\mathrm{corr}}\), and declared residual diagnostics
\State Generate teacher problems by solving the retained EOT/Sinkhorn route on sampled clouds from \(\mathcal D\)
\State Extract the latent teacher target to be predicted, e.g. dual/scaling coordinates
\State Fit the predictor \(S_\psi\) on this teacher data using the declared loss family
\State For a new weighted cloud \((X,w)\), compute the student prediction \(h_\psi=S_\psi(X,w;\varepsilon,N,d_x)\)
\State Translate \(h_\psi\) into the initial state of the retained corrective Sinkhorn solver
\State Run \(K_{\mathrm{corr}}\) corrective Sinkhorn updates to obtain \(\widehat\Pi_{\psi,\varepsilon,K_{\mathrm{corr}}}\)
\State Form the barycentric equal-weight output cloud \(\widetilde X_{\psi,\mathrm{corr}}=B(\widehat\Pi_{\psi,\varepsilon,K_{\mathrm{corr}}})\)
\State Record row/column residuals, teacher-student discrepancy, and any downstream scalar checks
\end{algorithmic}
\end{algorithm}

\paragraph{What the student is allowed to change.}
Only the initialization of the retained teacher solver, or another explicitly
named latent teacher coordinate.

\paragraph{What the teacher still protects.}
The declared transport constraints, the EOT objective, and the corrected
barycentric output after refinement.

\section{Training objective and deployment contract}
\label{sec:bf-learned-ot-training-objective}

Let \(h_\tau\) denote the teacher latent target and let
\(\mathcal D(N,d_x,\varepsilon,m,\eta)\) be the training distribution over
problem instances.  A simple supervised teacher-student fitting problem is
\begin{equation}
\label{eq:bf-learned-ot-training-objective}
  \psi^\star
  \in
  \arg\min_\psi
  \mathbb E_{(X,w,\varepsilon,m)\sim\mathcal D}
  \left[
    \mathcal L\{S_\psi(X,w;\varepsilon,N,d_x),h_\tau(X,w,\varepsilon)
    \}
  \right].
\end{equation}
In practice, the deployment claim is not about this loss alone.  It is about
whether the corrected route preserves the teacher output well enough on held-out
and stress-test clouds.

\section{Permutation symmetry}
\label{sec:bf-learned-ot-symmetry}

Particle labels are arbitrary.  For a permutation matrix
\(P\in\mathbb R^{N\times N}\), the predictor should at least satisfy a declared
permutation symmetry.  In the output-cloud case the natural condition is
\begin{equation}
\label{eq:bf-learned-ot-equivariance}
  S_\psi(PX,Pw;\varepsilon,N,d_x)
  =
  P\,S_\psi(X,w;\varepsilon,N,d_x)
  \quad
  \text{or the corresponding latent-state analogue.}
\end{equation}
This is a necessary condition, not a target-correctness theorem.

\section{Residual hierarchy}
\label{sec:bf-learned-ot-residuals}

The first deployment residual is the teacher-student discrepancy after correction:
\begin{equation}
\label{eq:bf-learned-ot-residual}
  R_{\psi,\tau}(X,w,\varepsilon)
  =
  B(\widehat\Pi_{\psi,\varepsilon,K_{\mathrm{corr}}})
  -
  B(\Pi_\tau(X,w,\varepsilon)).
\end{equation}
This residual does not by itself control filtering or HMC target error.  It only
measures how well the corrected student reproduces the chosen teacher.

\section{Filtering-loop insertion point}
\label{sec:bf-learned-ot-loop}

Inside the filtering loop, the retained-teacher neural route sits exactly where
Chapter~\ref{ch:bf-dpf-entropic-ot-sinkhorn} placed the finite EOT/Sinkhorn
layer:
\begin{equation}
\label{eq:bf-learned-ot-loop}
  \{x_i,w_i\}_{i=1}^N
  \xrightarrow{\mathrm{student\ warm\ start + retained\ Sinkhorn}}
  \{\tilde x^{(j)},1/N\}_{j=1}^N.
\end{equation}
It accelerates the equal-weighting layer after the weighted cloud has already been
formed.  It does not redefine the earlier particle-weighting rule itself.

\section{Verification notes}
\label{sec:bf-learned-ot-verification}

A practical implementation should still check a few local quantities carefully:
\begin{enumerate}[label=(\arabic*)]
  \item teacher generation is declared and reproducible,
  \item the predicted latent teacher object is named explicitly,
  \item the corrected route satisfies the row/column residual contract,
  \item teacher-student cloud discrepancy remains below the declared threshold on
  the tested envelope,
  \item the downstream scalar is checked whenever the learned route is embedded in
  a filtering likelihood or HMC target.
\end{enumerate}

These checks are mathematical and numerical, not bureaucratic.  Their purpose is
only to ensure that the reader can tell whether the implementation is still
computing the declared teacher-student object or has drifted to a different one.

\section{What this chapter does and does not claim}
\label{sec:bf-learned-ot-boundary}

This chapter teaches the main implementation-bearing neural route in the OT block:
retain the EOT/Sinkhorn teacher and learn how to start it better.  It does
\emph{not} claim:
\begin{itemize}
  \item that the student replaces the teacher object,
  \item that small teacher-student residuals automatically imply posterior or HMC
  correctness,
  \item that broader neural OT families are unimportant, only that they answer a
  different implementation question.
\end{itemize}
Its role is narrower: fix one teacher-student acceleration route in a form that a
reader can nearly implement directly.

\FloatBarrier
